\documentclass[11pt,a4paper]{ctexart}
\usepackage[margin=2.4cm]{geometry}
\usepackage{amsmath,amssymb}
\usepackage{fancyhdr}
\usepackage{enumitem}
\linespread{1.18}
\setlength{\headheight}{15pt}
\setlength{\parskip}{0pt}

\pagestyle{fancy}
\fancyhf{}
\fancyhead[L]{\small Team 2425397}
\fancyhead[R]{\small \thepage}
\renewcommand{\headrulewidth}{0.4pt}
\fancypagestyle{plain}{\fancyhf{}\renewcommand{\headrulewidth}{0pt}}

\ctexset{
  section/format = \Large\bfseries\raggedright,
  subsection/format = \large\bfseries\raggedright,
  subsubsection/format = \normalsize\bfseries\raggedright
}
\setcounter{secnumdepth}{3}

\begin{document}

%====================== 摘要页 ======================
\thispagestyle{empty}
\begin{center}
{\small 2024 MCM/ICM\quad 选题 A\quad 队伍控制号 2425397\quad 摘要页（Summary Sheet）}\\[18pt]
{\LARGE\bfseries 七鳃鳗种群分析：基于年龄分组的离散种群动力学模型}\\[8pt]
{\large （Population Analysis of Lampreys: A Discrete Population Dynamics Model Based on Age Grouping）}\\[10pt]
\end{center}

\begin{center}
{\large\bfseries 摘\quad 要}
\end{center}

\vspace{6pt}
本文对七鳃鳗的生态系统进行建模，并分析七鳃鳗性别比的变化对种群与生态系统的影响。首先，我们建立一个带年龄分组的离散种群动力学模型。依据七鳃鳗的生物学特性，我们把七鳃鳗种群划分为七个年龄组，并考虑它们所栖息的生态系统以及幼体的饵料、幼体的竞争者、成体的饵料、成体的竞争者与人工渔业。该模型考虑了性别比随外部环境而变化这一事实。

针对任务一，我们关注七鳃鳗性别比的变化对更大范围生态系统的影响。离散种群动力学模型的结果反映出，七鳃鳗可变的性别比有利于其种群规模的稳定，使其在生态系统中占据优势。通过比较七鳃鳗入侵某一生态系统前后种群规模的变化，我们得出结论：七鳃鳗作为入侵物种所具有的可变性别比对该生态系统的本地种群不利；但与此同时，七鳃鳗的幼体能够净化水质，并促进底栖环境中的养分循环。

针对任务二，我们应用模型，重点关注七鳃鳗可变性别比、变态、一次性繁殖与洄游这些生物学特性对其自身种群的利与弊。通过离散种群动力学模型计算环境资源变化与外部环境突变的情形，我们得到：具有变态与可变性别比的七鳃鳗种群更为稳定，具有变态的七鳃鳗种群更能实现对环境资源的充分利用。通过模拟人类影响对七鳃鳗洄游产卵过程的作用以及某一年份产卵亲鱼数量的急剧下降，我们得出结论：七鳃鳗的洄游特性与一次性繁殖的弊端在于，洄游过程可能存在危险，七鳃鳗的幼体对栖息地具有依赖性，而某一世代存活率偏低会使整个种群面临风险。

针对任务三，我们密切关注七鳃鳗可变性别比对生态系统稳定性的影响。我们以生态系统的抵抗稳定性、恢复稳定性与可持续性作为评价指标。离散种群动力学模型的数值模拟表明：当性别比可变的七鳃鳗数量与浮游生物种群数量发生突变时，生态系统的波动程度分别为 $23.28\%$ 与 $33.19\%$，恢复时间分别为 18 年与 28 年。与七鳃鳗性别比不可变时的数值模拟结果相比，具有可变性别比的七鳃鳗种群使生态系统具有更强的抵抗稳定性与恢复稳定性。在引入人工干预之后，含有可变性别比七鳃鳗种群的生态系统是可持续的。

针对任务四，我们主要用模型分析七鳃鳗的可变性别比对其他物种的影响。我们在离散种群动力学模型中加入寄生物种群，模型结果表明，七鳃鳗的可变性别比有利于其他寄生生物与种群。总之，七鳃鳗性别比的变化有助于在其面对恶劣环境时保证种群数量的稳定，这进而有助于更大范围生态系统的稳定与可持续。

\vspace{6pt}
\noindent\textbf{关键词：}离散种群动力学建模；七鳃鳗性别比；七鳃鳗变态；生态系统稳定性

\newpage
\pagestyle{fancy}

%%====================== 正文 ======================
\section{引言}

\subsection{问题背景}

在自然界中，不同物种的性别比各不相同。大多数物种的性别比为 $1:1$，但海七鳃鳗的性别比会随外部环境的变化而改变。我们把这种通过改变物种性别比以适应环境需求的现象称为适应性性别比变化。

七鳃鳗主要生活在大多数温带地区的沿海与淡水水域。作为若干湖泊生境中的寄生物，以及作为食物链中的一员，七鳃鳗对生态系统有显著影响，为斯堪的纳维亚、波罗的海以及北美太平洋西北部等地区提供食物来源。

七鳃鳗幼体阶段的发育速度决定其性别比，而其生长速度受食物可获得性的影响。当食物供给不足时，种群的性别以雄性为主，比例约为 $78\%$；当食物供给充足时，雄性比例约为 $56\%$，雌雄比例几乎相等。

\subsection{问题重述}

针对七鳃鳗的性别比及其对局部条件的依赖性问题，我们建立合理的数学模型，以深入认识生态系统中产生的相互作用，并回答以下问题：

\begin{itemize}
\item 当七鳃鳗的性别比与其种群规模相关时，分析其对七鳃鳗所处生态系统的影响。
\item 利用相应建立的模型，分析并论证七鳃鳗种群的优点与弱点。
\item 考虑七鳃鳗的性别比随外部环境而变化的情形，分析其对生态系统稳定性的影响。
\item 基于七鳃鳗种群性别比的特殊性，说明七鳃鳗能否为生态系统中的其他物种（例如寄生物）带来优势。
\end{itemize}

\subsection{我们的工作}

本文的工作流程可概括为：以年龄分组为基础建立离散种群动力学模型，依次完成任务一至任务四的建模与分析，随后对模型进行灵敏度检验，并总结模型的优缺点。

\section{假设与符号说明}

\subsection{模型假设}

\begin{itemize}
\item \textbf{假设 1：}我们假设所有七鳃鳗均在第七年繁殖，且亲体在繁殖结束时死亡。

理由：七鳃鳗是一次性繁殖生物，寿命为七年，因此为便于计算七鳃鳗繁殖所产生的新生幼体数量，我们假设它们全部在第七年繁殖。

\item \textbf{假设 2：}我们假设七鳃鳗的性别转变在其进入第五年的瞬间完成。

理由：在模型中我们已按年龄对七鳃鳗进行分组，并以年为间隔计算各年龄七鳃鳗的种群规模，因此为便于计算数量，我们认为七鳃鳗改变性别的行为是瞬时完成的。

\item \textbf{假设 3：}我们假设七鳃鳗没有天敌，即我们不把七鳃鳗视为饵料，其主要威胁来自人类活动。

理由：七鳃鳗幼体生活在相对隐蔽的环境中，很大程度上免于常规捕食，而七鳃鳗成体因其体型与习性，受天敌威胁相对较少。因此我们不引入七鳃鳗的捕食者种群，而是加入人工干预项来刻画人类活动对七鳃鳗种群的影响。

\item \textbf{假设 4：}我们假设七鳃鳗的寄生物种群仅寄生于七鳃鳗的成熟体。

理由：关于七鳃鳗寄生物的现有研究表明，超过 $80\%$ 的寄生物见于七鳃鳗的成熟体种群，仅 $11\%$ 的寄生物存在于七鳃鳗的幼体种群，因此为简化模型，我们主要考虑寄生物对七鳃鳗成熟体种群的影响。
\end{itemize}

\subsection{符号说明}

本文所用主要符号如下：$M_t$ 表示雌性比例；$\beta_m$ 表示幼体的自然死亡率；$\beta_B$ 表示缺乏饵料条件下的自然死亡率；$\alpha$ 表示繁殖系数；$r_x$ 表示自然增长率。其余符号在文中首次出现时另行说明。

\section{七鳃鳗生态系统模型}

\subsection{模型准备：七鳃鳗种群的年龄分组}

\subsubsection{七鳃鳗的年龄构成及其子群}

根据资料收集 [1]，七鳃鳗的寿命为 7 年，前 4 年为无性别差异的幼体，以浮游生物为食；到第 5 年，它们开始发生形态转变，出现雌雄之分，性别比主要受食物供给的影响。因此，七鳃鳗幼体期的食物量能够影响形态转变期的性别比，进而影响七鳃鳗的种群规模，其影响关系如图 3 所示。

本文把七鳃鳗种群的年龄划分为七个阶段，每一段向上取整（即七鳃鳗的年龄从 1 岁开始），如图 4 所示，其中 1—4 岁的七鳃鳗为无性别差异的幼体，5—7 岁的七鳃鳗为已发生性别差异的成体。

\subsubsection{七鳃鳗生态系统}

七鳃鳗的幼体主要捕食浮游动物，而七鳃鳗的成熟体则寄生鱼类。在考虑人工捕捞之后，我们把七鳃鳗所处的生态系统划分为三个部分，其关系包括竞争、寄生等。生态系统的关系图如图 5 所示。

\subsection{基于年龄分组的离散种群动力学模型}

\subsubsection{七鳃鳗幼体种群的生态系统建模}

\subsubsection*{1). 七鳃鳗的幼体种群}

根据我们对七鳃鳗的年龄分组及其物种特性，幼体不能繁殖。当环境中没有食物饵料时，幼体的自然死亡率为
\begin{equation}
\frac{X_{i+1,t+1}-X_{i,t}}{X_{i,t}}=-1.
\end{equation}

当环境中存在食物饵料时，幼体只会部分死亡，此时的死亡率为
\begin{equation}
\frac{X_{i+1,t+1}-X_{i,t}}{X_{i,t}}=-1+\frac{1-\beta_m}{K_x}x_t,
\end{equation}
其中 $x_t$ 为幼体饵料的数量，$\beta_m$ 为饵料充足时幼体的自然死亡率，$K_x$ 为 $x_t$ 的环境容纳量。考虑到与幼体共享同一食物饵料的竞争者种群 $B_t$ 以及不同年龄幼体之间的竞争，幼体的自然增长率为
\begin{equation}
\frac{X_{i+1,t+1}-X_{i,t}}{X_{i,t}}=-1+\frac{1-\beta_m}{K_x}x_t-\sum_{i=1}^{4}\tilde{k}_iX_{i,t}-k_tB_t,
\end{equation}
其中 $\tilde{k}_i$ 为不同年龄幼体之间的相互竞争系数，$k_t$ 为竞争者的竞争系数。

根据我们的假设，七鳃鳗仅在 7 岁时繁殖，因此幼体种群中新生幼体的数量为
\begin{equation}
X_{1,t+1}=\alpha M_tX_{7,t}.
\end{equation}

根据相关研究，种群丰度与种群中雄性个体的比例呈正相关 [2]，
\begin{equation}
M_t=\frac{\varepsilon_1+\varepsilon_2\sum\limits_{j=5}^{7}X_{j,t}}{\varepsilon_3+\varepsilon_4\sum\limits_{j=5}^{7}X_{j,t}},
\end{equation}
其中 $\alpha$ 为繁殖系数，$\varepsilon_i\ (i=1,2,3,4)$ 为非线性函数参数。

\subsubsection*{2). 七鳃鳗的竞争者种群}

在缺乏饵料时，七鳃鳗幼体竞争种群的自然增长率为
\begin{equation}
\frac{B_{t+1}-B_t}{B_t}=-\beta_B.
\end{equation}

当存在饵料时，其自然增长率为
\begin{equation}
\frac{B_{t+1}-B_t}{B_t}=-\beta_B+\gamma_B+\frac{\beta_B}{K_x}x_t.
\end{equation}

当考虑物种之间的竞争时，上式变为
\begin{equation}
\frac{B_{t+1}-B_t}{B_t}=-\beta_B+\gamma_B+\frac{\beta_B}{K_x}x_t-\sum_{i=1}^{4}k_iX_{i,t}-\tilde{k}_tB_t,
\end{equation}
其中 $\beta_B$ 为缺乏饵料条件下的自然死亡率，$\gamma_B$ 为饵料充足条件下的自然增长率，$k_i$ 为七鳃鳗幼体对该种群的竞争系数，$\tilde{k}_t$ 为种内竞争系数。

\subsubsection*{3). 七鳃鳗的饵料种群}

对于饵料种群，当其不存在捕食者时，其种群规模符合 Logistic 增长：
\begin{equation}
\frac{x_{t+1}-x_t}{x_t}=r_x\left(1-\frac{x_t}{K_x}\right).
\end{equation}

当考虑其捕食者存在时，则有
\begin{equation}
\frac{x_{t+1}-x_t}{x_t}=r_x\left(1-\frac{x_t}{K_x}\right)-\sum_{i=1}^{4}k_iX_{i,t}-k_tB_t,
\end{equation}
其中 $r_x$ 为资源充足时的自然增长率。

\subsubsection{七鳃鳗成体种群的生态系统建模}

\subsubsection*{1). 七鳃鳗的成体种群}

七鳃鳗由幼体转变为成体时的自然增长率可表示为
\begin{equation}
\frac{X_{5,t+1}-X_{4,t}}{X_{4,t}}=-1+\frac{1-\beta_m}{K_x}x_t-\sum_{i=1}^{4}\tilde{k}_iX_{i,t}-k_tB_t.
\end{equation}

考虑实际情况，在存在饵料与竞争时其自然增长率与幼体情形类似，同理可得关系式
\begin{equation}
\frac{X_{j+1,t+1}-X_{j,t}}{X_{j,t}}=-1+\frac{1-\beta_n}{K_y}y_t-\sum_{j=5}^{7}\tilde{k}_jX_{j,t}-k_LC_t,
\end{equation}
其中 $y_t$ 为七鳃鳗成体的饵料种群数量，$C_t$ 为七鳃鳗成体的竞争者种群数量，$\beta_n$ 为饵料充足时的自然死亡率，$\tilde{k}_j$ 为种内竞争系数，$k_L$ 为竞争者对成体的竞争系数。

\subsubsection*{2). 七鳃鳗的竞争者种群}

成体由幼体发育而来，与幼体种群类似，成体竞争者种群规模的变化关系为
\begin{equation}
\frac{C_{t+1}-C_t}{C_t}=-\omega_y+\eta_y+\frac{\omega_y}{K_y}y_t-\sum_{j=5}^{7}k_jX_{j,t}-\tilde{k}_LC_t.
\end{equation}

饵料种群规模的变化关系为
\begin{equation}
\frac{y_{t+1}-y_t}{y_t}=r_y\left(1-\frac{y_t}{K_y}\right)-\sum_{j=5}^{7}k_jX_{j,t}-k_LC_t,
\end{equation}
其中 $K_y$ 为饵料种群 $y_t$ 的环境容纳量。

\subsubsection{外部人工干预模型}

考虑到人工干预对七鳃鳗种群的长期压力，假设年干预强度为 $F$，则对式 (3)、(11) 与 (17) 有
\begin{equation}
\left\{
\begin{array}{l}
\dfrac{X_{i+1,t+1}-X_{i,t}}{X_{i,t}}=-1+\dfrac{1-\beta_m}{K_x}x_t-\sum\limits_{i=1}^{4}\tilde{k}_iX_{i,t}-k_tB_t-F\quad(i=1,2,3),\\[10pt]
\dfrac{X_{5,t+1}-X_{4,t}}{X_{4,t}}=-1+\dfrac{1-\beta_m}{K_x}x_t-\sum\limits_{i=1}^{4}\tilde{k}_iX_{i,t}-k_tB_t-F,\\[10pt]
\dfrac{X_{j+1,t+1}-X_{j,t}}{X_{j,t}}=-1+\dfrac{1-\beta_n}{K_y}y_t-\sum\limits_{j=5}^{7}\tilde{k}_jX_{j,t}-k_LC_t-F\quad(j=5,6).
\end{array}
\right.
\end{equation}
这更好地表达了七鳃鳗幼体与成体种群的自然的增长率。

\subsection{模型概述}

综上所述，若幼体与成体种群的数量分别记为 $X_i\ (i=1,2,3,4)$ 与 $X_j\ (j=5,6,7)$，饵料数量为 $x_t$ 与 $y_t$，竞争者为 $B_t$ 与 $C_t$，则新生幼体的数量为
\[
X_{1,t+1}=\alpha M_tX_{7,t},
\]
其中 $\alpha$ 表示繁殖系数，$M_t$ 表示雌性比例，而雌性比例为
\[
M_t=\frac{\varepsilon_1+\varepsilon_2\sum\limits_{j=5}^{7}X_{j,t}}{\varepsilon_3+\varepsilon_4\sum\limits_{j=5}^{7}X_{j,t}}.
\]

当考虑人工干预（即 $F\neq 0$）时，种群规模之间的变化关系如表 2 与表 3 所示，其中 $\tilde{k}_i$、$\tilde{k}_j$、$\tilde{k}_t$ 与 $\tilde{k}_L$ 为种内竞争系数。表 2 给出成体部分（成体、竞争者、饵料）的数量变化关系，表 3 给出幼体部分（幼体、竞争者、饵料）的数量变化关系，相应的关系式依次为：
\[
\frac{X_{5,t+1}-X_{4,t}}{X_{4,t}}=-1+\frac{1-\beta_m}{K_x}x_t-\sum_{i=1}^{4}\tilde{k}_iX_{i,t}-k_tB_t-F,
\]
\[
\frac{X_{j+1,t+1}-X_{j,t}}{X_{j,t}}=-1+\frac{1-\beta_n}{K_y}y_t-\sum_{j=5}^{7}\tilde{k}_jX_{j,t}-k_LC_t-F,
\]
\[
\frac{C_{t+1}-C_t}{C_t}=-\omega_y+\eta_y+\frac{\omega_y}{K_y}y_t-\sum_{j=5}^{7}k_jX_{j,t}-\tilde{k}_LC_t,
\]
\[
\frac{y_{t+1}-y_t}{y_t}=r_y\left(1-\frac{y_t}{K_y}\right)-\sum_{j=5}^{7}k_jX_{j,t}-k_LC_t,
\]
\[
\frac{X_{i+1,t+1}-X_{i,t}}{X_{i,t}}=-1+\frac{1-\beta_m}{K_x}x_t-\sum_{i=1}^{4}\tilde{k}_iX_{i,t}-k_tB_t-F,
\]
\[
\frac{B_{t+1}-B_t}{B_t}=-\beta_B+\gamma_B+\frac{\beta_B}{K_x}x_t-\sum_{i=1}^{4}k_iX_{i,t}-\tilde{k}_tB_t,
\]
\[
\frac{x_{t+1}-x_t}{x_t}=r_x\left(1-\frac{x_t}{K_x}\right)-\sum_{i=1}^{4}k_iX_{i,t}-k_tB_t.
\]

\section{任务一：对更大范围生态系统的影响}

\subsection{种群性别比变化的分析}

根据前述基于年龄分组的七鳃鳗生态系统模型，我们可以得到图 6 所示的结果，其初始取值如表 4 所示。模型初始参数取值为：$X_{1,1}=0.4$、$X_{2,1}=0.35$、$X_{3,1}=0.3$、$X_{4,1}=0.25$、$X_{5,1}=0.2$、$X_{6,1}=0.15$、$X_{7,1}=0.1$、$B_1=1$；$x_1=1$、$y_t=1$、$C_1=1$、$M_1=0.3$、$T=100$、$\alpha=12$、$\beta_m=0.01$、$K_x=1$；$\tilde{k}_i=0.001$、$k_i=0.1$、$\tilde{k}_t=0.001$、$k_t=0.01$、$\beta_B=0.1$、$\gamma_y=0.1499$、$r_x=0.5$、$\beta_n=0.01$；$K_y=1$、$\tilde{k}_j=0.001$、$k_j=0.1$、$\tilde{k}_L=0.001$、$k_L=0.01$、$\omega_y=0.1$、$\eta_y=0.0471$、$r_y=0.5$。

在系统达到稳定之后，根据图 6(a)，我们得到七鳃鳗的种群规模在 $\pm 4\%$ 的范围内形成周期为 7 年的波动。图 6(b) 给出了七鳃鳗幼体与成体种群规模的变化，其同样以 7 年为周期波动。因此我们对数据进行平滑，并按如下公式进行变换：
\begin{equation}
\Delta_{i,t}=\frac{1}{7}\sum_{\tau=-3}^{3}\Delta_{i,t+\tau},
\end{equation}
其中 $\Delta_{i,t}$ 表示各类种群的数量，由此可得到平滑后的曲线，如图 7 所示。

从图 7 可以看出，生态系统在第 24 年达到平衡，七鳃鳗幼体的种群数量多于成熟体。这是由于七鳃鳗成体的产卵量远大于同期死亡的亲体数量。此外，七鳃鳗幼体随生长而数量减少，成为成体的数量将远小于其产卵能力。由于七鳃鳗幼体数量众多，它们在与竞争者的竞争中处于优势。另一方面，七鳃鳗成体数量较少，这使它们在竞争中处于劣势。

假设雌雄比例固定且各占 $50\%$，在其他条件均与图 7 相同的情况下，我们绘制出图 8 所示的图像。与七鳃鳗性别可变的情形相比，在性别固定的情形下，七鳃鳗的种群规模会更大，但其饵料种群的数量将减少，这将导致分配给每个七鳃鳗个体的环境资源量更少，种群的生活质量下降。此外，固定的七鳃鳗性别比将导致生态系统中其他生物的数量减少甚至灭绝，这对生态系统的多样性不利。

\subsection{对生态系统的影响}

\subsubsection*{1). 作为入侵物种危害生态系统中的本地种群}

为模拟七鳃鳗作为外来物种入侵五大湖，我们首先构建一个没有七鳃鳗的生态系统，如图 9(a) 所示。在七鳃鳗入侵图 9(a) 中的系统之前，系统内各类种群处于稳定状态。若七鳃鳗从当前稳定生态系统的状态开始入侵，则可得到图 9(b) 所示的图像，其中以图 9(a) 中各类种群数量的稳定值作为图 9(b) 中各类种群数量的初值。

入侵生态系统之后，七鳃鳗没有天敌，因而能在短时间内大量繁殖，并最终因环境资源量而使其种群趋于稳定。在此过程中，本地种群显著衰退，甚至可能导致某个种群灭绝。以七鳃鳗入侵北美五大湖为例，鱼类种群规模的变化与图 9(b) 类似。七鳃鳗入侵期间出现的大幅衰退给当地渔业经济造成了巨大损失 [3]。

\subsubsection*{2). 净化水质}

从图 7(a) 曲线的趋势可以发现，七鳃鳗幼体的数量在总量中占显著更大的比例，而大量七鳃鳗幼体摄食水生环境中的微生物与浮游生物，可起到净化水质的作用 [4][5]。

\subsubsection*{3). 影响底栖生态系统}

七鳃鳗的幼体生活在淡水河流与溪流的底部，主要滤食有机碎屑与微生物。它们通过挖掘沉积物构建生境 [6]，这一活动增加了底部沉积物的氧化，促进了底栖环境中的养分循环。

\subsection{结果分析}

我们对式 (5) 进行验证，结果均以未平滑数据的形式展示于图 10。

从图 10(a) 可以看出，七鳃鳗的雌性比例与成熟体种群数量呈负相关，当成熟体种群数量上升时七鳃鳗的雌性比例下降，这与文献 [2] 中的研究相同。从图 10(b) 可以看出，七鳃鳗的性别比与环境资源量具有相同的变化趋势，当环境资源量增加时七鳃鳗的雌性比例上升，当资源量减少时七鳃鳗的雌性比例下降，这与问题背景的描述相符。这表明我们对七鳃鳗种群性别比变化方式的描述是合理的。

由于我们对波动数据进行了平滑，因此需要验证其合理性。

图 11 中的闭环曲线表明，七鳃鳗幼体与浮游生物、七鳃鳗成熟体与鱼类的种群规模均呈周期性波动，反映出系统是稳定的，因此我们对原始数据取 7 年平均值是合理的。
\section{任务二：种群的优点与弱点}

\subsection{种群的优点}

\subsubsection*{1). 变态的优点：耐受恶劣环境}

为模拟环境的变化，我们改变七鳃鳗种群可获得的饵料量，并改变浮游生物与鱼类的增长率以及环境容纳量。

我们减少环境资源量，并把七鳃鳗的幼体与成体分别与其饵料和竞争者一起绘制于图 12(a) 与图 13(a)。随后我们增加环境资源量，分别绘制出图 12(b) 与图 13(b)。

根据图 12(a) 与图 12(b)，我们可以看到七鳃鳗幼体受环境资源量变化的影响更大，在食物饵料量下降之后幼体种群持续衰退，甚至可能灭绝。相比之下，图 13 表明七鳃鳗的成熟个体在食物饵料量减少后种群规模会下降，但最终将达到相对稳定的状态。因此，七鳃鳗的变态能够提高其对恶劣环境的抵抗力，七鳃鳗可以通过变态有效提高种群的适应能力，从而保证种群的正常繁殖。

在饵料数量上升之后，成体种群与幼体种群均持续增加，这表明七鳃鳗在食物丰富的环境中更具竞争力，也更易于繁殖。

\subsubsection*{2). 变态的优点：充分利用环境资源}

为研究能否变态对七鳃鳗种群的影响，我们假设七鳃鳗的幼体不变态为成体，构建仅含七鳃鳗幼体、幼体竞争者与浮游生物的生态模型以得到图 14(a)，然后比较能否变态这两种生态模型中七鳃鳗的数量与浮游生物的数量，得到图 14(b)。

比较图 14(a) 与图 14(b)，我们得出结论：不发生性别转变的七鳃鳗其种群规模小于发生性别转变的情形，说明七鳃鳗的性别转变有利于种群繁殖。我们分析其原因是，七鳃鳗在性别状态转变前后具有不同的摄食饵料，这有效降低了种内竞争，实现了对环境资源的充分利用。

\subsubsection*{3). 可变性别比的优点：具有抵抗突变的能力}

为模拟浮游生物暴发事件，我们在第 75 天使浮游生物大幅增加，得到图 15。

根据图 15，我们观察到生态系统内浮游生物种群的显著增加会引发环境资源的变化，进而诱导七鳃鳗种群性别比的改变。这种适应性机制使七鳃鳗种群先随环境变化而扩张，随后收缩，并最终稳定在与扰动前规模几乎相当的水平上。由此我们推断，以灵活的性别比为特征的七鳃鳗种群对环境波动表现出更强的恢复力。

\subsection{种群的弱点}

\subsubsection*{1). 洄游的弊端：人类干预可对种群造成更大的抑制}

为模拟人类干预对七鳃鳗物种的影响，我们减少第七年洄游产卵的七鳃鳗数量，得到图 16。

人工干预减少了洄游产卵的七鳃鳗数量，进而减少了七鳃鳗幼体的数量，七鳃鳗种群受到更强的抑制。图 16 反映了在人工干预作用下七鳃鳗洄游产卵的下降。

\subsubsection*{2). 受栖息环境的高度影响}

根据图 16 我们发现，七鳃鳗幼体对环境资源量的变化非常敏感，栖息地退化会极大地影响七鳃鳗幼体的种群规模，从而导致整个种群的衰退。

\subsubsection*{3). 一次性繁殖的风险}

根据七鳃鳗的种群生物学，七鳃鳗只繁殖一次，随后即死亡。我们模拟某一世代七鳃鳗遭受人工干预或环境影响以致产卵存活率过低的情形，将第 50 年七鳃鳗的产卵数量减半，得到图 17。当某一世代七鳃鳗因人工干预或环境影响导致产卵存活率过低时，七鳃鳗种群显著减少，这极大地威胁种群的恢复能力并导致种群规模衰退，使种群陷入生存危机。

\section{任务三：性别比对生态系统稳定性的影响}

\subsection{生态系统稳定性的分类}

我们把生态系统稳定性的评价指标分为生态系统的抵抗稳定性、恢复稳定性与可持续性。我们在模型中引入扰动量以模拟突发事件对生态系统中种群的影响，并引入人工干预项以模拟年度人类活动对七鳃鳗种群的影响。根据我们的模拟，七鳃鳗种群性别比的变化对生态系统稳定性的影响如下。

\subsection{生态系统的抵抗稳定性与恢复稳定性}

\subsubsection*{1). 七鳃鳗种群规模的突变}

我们在模型中引入脉冲扰动项 $E$，假设人类在第 $t_0$ 年对七鳃鳗的捕捞导致种群规模显著下降 $E$。随后我们分别分析性别比可变与不可变的七鳃鳗所处生态系统的稳定性。

图 18 中性别比可变与不可变情形的数据如表 5 所示，其中
\begin{equation}
t_{\Delta}=t_2-t_1,
\end{equation}
其中 $t_1$ 为扰动的起始时间，$t_2$ 为恢复的结束时间。此外，波动程度为各类群体波动之和
\begin{equation}
\phi=\sum\frac{|\Delta_{i,t+\tau}-\Delta_{i,t}|}{\Delta_{i,t}},
\end{equation}
其中 $\tau$ 为最大波动所经历的时间，$\phi$ 为波动百分比。

图 18 中性别比可变与不可变情形的具体数据为：性别比可变时，恢复时长为 18 年，七鳃鳗、幼体、竞争者与浮游生物的波动程度分别为 $23.28\%$、$8.46\%$、$12.52\%$ 与 $2.30\%$；性别比不可变时，恢复时长为 22 年，四者的波动程度分别为 $30.41\%$、$8.72\%$、$18.06\%$ 与 $3.63\%$。

\subsubsection*{2). 浮游生物种群的突变}

为研究浮游生物种群的突变，我们引入脉冲扰动项 $F$，并假设该生态系统中的浮游生物种群在第 $t_1$ 年发生突变。

图 19 中性别比可变与不可变情形的数据由式 (17) 与式 (18) 计算得到，具体为：性别比可变时，恢复时长为 28 年，七鳃鳗、幼体、竞争者、浮游生物与成体竞争者的波动程度分别为 $33.19\%$、$18.46\%$、$2.74\%$、$5.24\%$ 与 $6.75\%$；性别比不可变时，恢复时长为 38 年，上述各类的波动程度分别为 $49.75\%$、$18.84\%$、$20.85\%$、$5.40\%$ 与 $4.66\%$。

基于上述图像与数据，我们发现性别比可变的七鳃鳗所处生态系统的抵抗稳定性与恢复稳定性均强于性别比不可变的七鳃鳗所处生态系统。因此，当七鳃鳗的性别比可变时，其所处生态系统的稳定性更好。

\subsection{生态系统的可持续性}

为研究七鳃鳗种群所处生态系统的可持续性，我们在该生态系统中七鳃鳗的种群数量 $X_j$ 上加入强度为 $F$ 的人工干预，以模拟该生态系统在长期承受来自人类的压力时各类种群数量的变化。

图 20 表明，当性别比可变且存在长期干预时，七鳃鳗种群会下降但趋于稳定。生态系统中的其他种群在经历同样轻微的变动时也保持稳定。这表明性别比可变的七鳃鳗所处生态系统是可持续的。

\subsection{小结}

综上所述，七鳃鳗种群性别比的变化使其所处生态系统更为稳定。无论是生态系统内种群规模的突变，还是对生态系统的长期压力，七鳃鳗种群都能通过性别比的变化而变得更加稳定。

\section{任务四：性别比变化对其他物种的影响}

\subsection{加入寄生物种群}

在研究七鳃鳗种群的寄生物时，超过 $80\%$ 的报告（包括生活史阶段）来源于七鳃鳗的成熟体，仅有 $11\%$ 的报告来源于七鳃鳗的幼体。这可能表明七鳃鳗种群的寄生物主要由七鳃鳗的成熟体所寄生 [7]。因此我们在模型中引入七鳃鳗成熟体种群的寄生物种群，并考虑其寄生对七鳃鳗成熟体种群的影响。

若寄生物种群的数量为 $P_t$，在无人工干预（即 $F=0$）时，七鳃鳗成熟体的自然增长率为
\begin{equation}
\frac{X_{j+1,t+1}-X_{j,t}}{X_{j,t}}=-1+\frac{1-\beta_n}{K_y}y_t-\sum_{j=5}^{7}\tilde{k}_jX_{j,t}-k_LC_t-DP_t,
\end{equation}
其中 $D$ 为感染寄生物的七鳃鳗所增加的死亡风险。而寄生物的自然增长率为
\begin{equation}
\frac{P_{t+1}-P_t}{P_t}=-d+\zeta+\frac{d}{k_m}\sum_{i=5}^{7}X_{j,t}-\tilde{k}_pP_t,
\end{equation}
其中 $P_t$ 为寄生物种群规模，$d$ 为寄生物在无宿主存在时的自然死亡率，$\zeta$ 为寄生物在宿主数量充足时的自然增长率，$\tilde{k}_p$ 为种内竞争系数，$k_m$ 为七鳃鳗成熟体的环境容纳量。

\subsection{模型求解与结果分析}

\subsubsection*{1). 从七鳃鳗总种群的角度}

图 21 给出了不区分幼体与成熟个体时加入寄生物之后种群规模变化的曲线。就七鳃鳗的总体种群规模而言，性别比可变与性别比不可变相比，前者的寄生物种群数量更多，表明当七鳃鳗的性别比可变时，对该生态系统中的寄生物种群有利，能够促进寄生物种群的繁殖。至于浮游生物，性别比可变的生态系统中浮游生物种群的数量也更多。两者均表明，性别比可变的特性同样有利于浮游生物种群。

\subsubsection*{2). 从七鳃鳗种群的幼体与成体数量角度}

从图 22 可以看出，寄生物主要寄生于七鳃鳗的成熟体，这导致成熟体种群规模下降，使繁殖数量减少，从而影响七鳃鳗幼体的数量。当七鳃鳗种群的性别比可变时，种群会随环境资源量的变化而改变雌性比例，使七鳃鳗种群的波动程度较小，寄生物种群的数量得以稳定增长。当七鳃鳗种群的性别比不可变时，受寄生物影响后新生种群的波动程度更大，这不利于寄生物种群数量的增加。与此同时，可变的性别比会减少幼体的数量，使浮游生物的数量上升，这有利于浮游生物的发展。

\section{模型灵敏度分析}

\subsection{模型参数的影响分析}

为分析模型参数的影响，我们关注在小范围内变动时雌性比例 $M_t$、繁殖系数 $\alpha$ 与人类干预强度 $F$ 对模型输出的影响。我们对模型中的每个参数施加 $\pm 5\%$ 的扰动，分析七鳃鳗的种群规模，以得到对模型输出影响最大的参数。

从图 23 可以看出，在 $t=100$ 时，繁殖系数 $\alpha$ 与雌性比例 $M_t$ 的变化引起的七鳃鳗种群规模波动分别达到 $17.56\%$ 与 $16.03\%$，而其他参数引起的波动均在 $3.34\%$ 以内，这表明七鳃鳗种群规模的变化对繁殖系数及其性别比非常敏感，$\alpha$ 与 $M_t$ 是模型中影响最大的参数。

\subsection{人工干预与寄生物种群对模型的影响}

为分析人工干预与寄生物种群对模型的影响，我们调整人工干预强度 $F$ 得到图 24(a)，并改变七鳃鳗被寄生物寄生后的死亡风险系数 $D$ 得到图 24(b)。

从图 24(a) 可以看出，七鳃鳗的种群规模随人工干预强度的增加而下降，因此人工干预是控制七鳃鳗种群规模的有效方法之一。从图 24(b) 可以看出，七鳃鳗种群随被寄生后死亡风险的上升而下降，这表明被不同寄生物种群寄生之后七鳃鳗种群的变化与寄生物所造成的损害程度有关，七鳃鳗种群可能没有显著变化，也可能出现显著下降。

\subsection{小结}

七鳃鳗种群数量的变化对其繁殖系数与性别比非常敏感，若需要控制其种群数量，可以通过控制其洄游产卵来控制其种群数量的变化，也可以通过改变环境资源量来控制其种群的性别比，从而控制其种群的繁殖。寄生物对七鳃鳗种群的影响与寄生物种群造成的损害程度有关，在特定七鳃鳗生物入侵时，可以考虑引入七鳃鳗的寄生物来抑制七鳃鳗种群的增长。

\section{模型的优缺点}

\textbf{优点：}

\begin{itemize}
\item 我们对七鳃鳗种群进行了年龄分组，能够有效考虑七鳃鳗在不同阶段种群规模的变化，并便于分析七鳃鳗的变态过程，使模型更符合实际。
\item 我们对七鳃鳗种群性别比变化的描述结合了相关文献中的专门研究，拟合函数具有科学有效性。
\end{itemize}

\textbf{缺点：}

\begin{itemize}
\item 在现实中，会有一些物种偶尔捕食七鳃鳗，而我们的模型未考虑七鳃鳗被捕获的情形，因此模型结果会与实际情况存在一定偏差。
\end{itemize}

\section*{参考文献}

\begin{itemize}
\item[{[1]}] Zhu Y, Jun L, Yue P, et al. Lampreys: an important model animal for biological evolution and disease research[J]. Genetics, 2020, 42 (09): 847-857.
\item[{[2]}] Hardisty M W. Sex Composition of Lamprey populations[J]. Nature, 1961, 191: 1116--1117.
\item[{[3]}] Johnson N S, Jubar A K, Keffer D A, et al. A case study of sea lamprey (Petromyzon marinus) control and ecology in a microcosm of the Great Lakes[J]. Journal of Great Lakes Research, 2021, 47: 492-505.
\item[{[4]}] Docker M F. Lampreys: Biology, Conservation and Control[M]. Springer, 2015.
\item[{[5]}] Boeker, C., Geist, J. Lampreys as ecosystem engineers: burrows of Eudontomyzon sp. and their impact on physical, chemical, and microbial properties in freshwater substrates[J]. Hydrobiologia, 2016, 777: 171--181.
\item[{[6]}] Daniel M. Weaver, Stephen M. Coghlan Jr., Joseph Zydlewski. Effects of sea lamprey substrate modification and carcass nutrients on macroinvertebrate assemblages in a small Atlantic coastal stream[J]. Journal of Freshwater Ecology, 2018, 33(1), 19-30.
\item[{[7]}] Megan A. Shavalier, Mohamed Faisal, Mary L. Moser, Thomas P. Loch. Parasites and microbial infections of lamprey (order Petromyzontiformes Berg 1940): A review of existing knowledge and recent studies[J]. Journal of Great Lakes Research, 2021, 47: 90-111.
\end{itemize}

\section*{人工智能工具使用报告}

\textbf{1. GPT-4 Turbo}

查询 1：哪些事件会导致浮游生物突然增加或减少？

输出：浮游生物种群数量可能因多种自然因素与人为因素而突然增加或减少。以下是可能导致浮游生物种群变化的一些事件：

导致浮游植物突然增加的事件：

\begin{itemize}
\item 光照增强：对于依赖光合作用的浮游植物（浮游藻类），春夏季较长的日照时数可促进其快速生长。
\item 营养盐输入：来自陆地的营养盐（如氮、磷）进入水体，尤其是通过农业径流或市政污水排放，可迅速促进浮游植物繁殖，形成“水华”。3. 水温变化：水温可能发生很大变化。
\item 水温变化：水温升高有利于某些浮游生物的生长速率，尤其在温带与寒带水体中。
\item 人工施肥：向养殖水体中人为添加化学肥料或有机肥料以促进浮游植物生长，进而为养殖鱼类或其他水生动物提供充足的食物。
\end{itemize}

导致浮游植物突然减少的事件：

\begin{itemize}
\item 捕食压力增加：浮游动物（如桡足类）种群的增加或鱼类的大量摄食，可迅速降低浮游植物种群数量。
\item 有毒排放：工业废水、农药及其他有毒化学品排入水体可导致浮游生物大量死亡。
\item 极端天气事件：洪水、干旱或寒潮等极端天气事件可通过改变水体的化学或物理条件间接影响浮游植物的生存。
\item 光照减弱：持续阴雨天气、浊度增加等因素降低水体中的光照量，影响依赖光合作用的浮游植物生长。
\item 富营养化：尽管营养盐的增加最初促进浮游植物生长，但长期而过度的富营养化可能导致水体缺氧，进而引发浮游动物及其他水生生物的大量死亡。
\end{itemize}

上述每一个事件都可以在生态系统模型中用脉冲项表示，以模拟其对浮游生物种群的瞬时效用。

\textbf{2. GPT-4 Turbo}

查询 2：如何判定生态系统的稳定性？

输出：生态系统的稳定性是指其在外界干扰与压力下保持结构与功能不变或迅速恢复的能力。判定生态系统稳定性涉及多个维度，包括生物多样性、生产力、抵抗力、恢复力与持续性。以下是评价生态系统稳定性的一些关键方面与方法：

\begin{itemize}
\item 物种多样性：物种丰富度与物种均匀度是生态系统稳定性的重要指标。高物种多样性通常意味着生态系统对外界干扰具有更强的抵御能力，因为多样的物种能够执行多样的生态功能，并增加系统的冗余度。
\item 初级生产力：初级生产力是指生态系统中的生产者（如植物）通过光合作用或化学合成将无机物转化为有机物的速率。高初级生产力表明生态系统具有较强的能量转换与物质循环能力，是生态系统稳定性最重要的指标之一。
\item 抵抗力：生态系统在面对干扰时保持其结构与功能的能力。
\item 恢复力：生态系统在受到干扰之后恢复到原状的速度与程度。可以通过评价生态系统在面对自然灾害、人类干预等干扰时的恢复情况来判定其稳定性。
\item 系统的持续性：生态系统在长期内维持其基本功能与结构、不因外界干扰而轻易崩溃的能力，表明其具有较高的稳定性。
\item 多个物种执行相同的生态功能：当多个物种能够执行相同的生态功能时，即使某些物种因干扰而减少或消失，生态系统的整体功能仍可由其他物种维持，这是衡量生态系统稳定性的一个重要方面。
\end{itemize}

判定方法：

\begin{itemize}
\item 长期生态学研究：通过对生态系统中物种组成、生物量与能量流动等指标的长期监测，评价其趋势与对干扰的响应。
\item 模型模拟：利用生态模型模拟生态系统在不同干扰下的动态变化，并预测其稳定性。
\item 对比分析：比较不同生态系统或同一生态系统在不同干扰条件下的响应，以评价其稳定性的差异。
\item 实验研究：通过设置受控实验（如人为增加或移除物种）来观察生态系统对干扰的响应，以评价其稳定性。
\end{itemize}

综上所述，判定生态系统的稳定性需要综合考虑多个生态指标与方法，只有通过长期观测与细致研究，才能全面理解生态系统在面对变化与干扰时的响应与适应能力。

\end{document}
